\section{Source coverage and archive reading guide}
\label{sec:source-map}

\subsection{The article and its immediate mathematical links}

The scope audit was made on 22 November 2026. The starting article and the
immediate linked pages listed below were inspected as a map of mathematical
topics \cite{sourcewiki}. The table records where this report develops those
topics. It is not a claim to have proved every sentence of every page, and
historical, biographical, and philosophical material is not part of the
theorem inventory. The mathematical proofs in this report are supplied
independently; the linked articles do not serve as unproved premises.

\begingroup\small
\begin{longtable}{@{}>{\raggedright\arraybackslash}p{.29\textwidth}>{\raggedright\arraybackslash}p{.49\textwidth}>{\raggedright\arraybackslash}p{.13\textwidth}@{}}
\toprule
Inspected topic & Coverage in this report & Sections\\
\midrule
\endfirsthead
\toprule
Inspected topic & Coverage in this report & Sections\\
\midrule
\endhead
\bottomrule
\endfoot
\href{https://en.wikipedia.org/wiki/Propositional_logic}{Propositional logic}
  (starting article) &
Syntax, rules, Boolean semantics, soundness, constructive completeness,
decidability, and distinctions between classical and intuitionistic logic.
& \ref{sec:syntax}, \ref{sec:nd}, \ref{sec:negative}\\[5pt]
\href{https://en.wikipedia.org/wiki/Propositional_formula}{Propositional formula}
& Unique parsing, finite support, subformulas, substitution, locality, and
primitive recursive coding of syntax.
& \ref{sec:syntax}\\[5pt]
\href{https://en.wikipedia.org/wiki/Truth_function}{Truth function};
\href{https://en.wikipedia.org/wiki/Truth_table}{truth table};
\href{https://en.wikipedia.org/wiki/Boolean_function}{Boolean function}
& Truth-table decision with witnesses; canonical forms; functional
completeness; monotone and affine fragments; algebraic normal form.
& \ref{sec:syntax}, \ref{ca-normal}\\[5pt]
\href{https://en.wikipedia.org/wiki/Natural_deduction}{Natural deduction}
& Explicit rule systems, proof checkers, structural transformations,
deduction, classical completeness, and constructive law templates.
& \ref{sec:nd}, \ref{sec:laws}\\[5pt]
\href{https://en.wikipedia.org/wiki/Hilbert_system}{Hilbert system};
\href{https://en.wikipedia.org/wiki/Combinatory_logic}{combinatory logic}
& Hilbert axioms and proof translations; the deduction theorem;
the roles of the $K$ and $S$ proof templates and abstraction.
& \ref{sec:systems}, \ref{sec:laws}\\[5pt]
\href{https://en.wikipedia.org/wiki/Method_of_analytic_tableaux}{Method of
analytic tableaux};
\href{https://en.wikipedia.org/wiki/Sequent_calculus}{sequent calculus}
& Bounded tableau expansion, open-branch countervaluations, analytic
classical completeness, and cut-free proof reconstruction.
& \ref{sec:systems}\\[5pt]
\href{https://en.wikipedia.org/wiki/Boolean_algebra_(structure)}{Boolean
algebra (structure)};
\href{https://en.wikipedia.org/wiki/Boolean_algebra_(logic)}{Boolean
algebra (logic)}
& Boolean identities; the algebra of truth-table bitsets; exact finite
cardinality, atoms, and the universal mapping property for finite targets.
& \ref{sec:laws}, \ref{ca-normal}\\[5pt]
\href{https://en.wikipedia.org/wiki/Boolean_satisfiability_problem}{Boolean
satisfiability problem};
\href{https://en.wikipedia.org/wiki/DPLL_algorithm}{DPLL algorithm}
& Witness-producing satisfiability procedures, branching bounds, clause
encodings, and polynomial verifier reductions to SAT and 3-SAT.
& \ref{ca-solvers}, \ref{sec:complexity}\\[5pt]
\href{https://en.wikipedia.org/wiki/Peirce's_law}{Peirce's law}
& Equivalence of classical axiom schemes over intuitionistic logic;
finite countermodels to isolated classical principles.
& \ref{sec:negative}, \ref{sec:ipc-saturation}\\[5pt]
\href{https://en.wikipedia.org/wiki/Intuitionistic_propositional_calculus}{Intuitionistic
propositional calculus} (redirecting to
\href{https://en.wikipedia.org/wiki/Intuitionistic_logic}{Intuitionistic logic})
& Finite Kripke semantics, completeness, decidability, finite Heyting
algebras, disjunction extraction, and the scope of negative translation.
& \ref{sec:ipc-saturation}, \ref{sec:negative}\\
\end{longtable}
\endgroup

The phrase ``immediate linked pages'' in this table applies to this inspected
list. It is not extended automatically to every further topic developed in
the report. In particular, the following are supporting developments chosen
because their proofs make the requested subject more complete and give
useful explicit arithmetic constructions.

\begin{center}\small
\begin{tabularx}{\textwidth}{@{}>{\raggedright\arraybackslash}p{.31\textwidth}>{\raggedright\arraybackslash}X@{}}
\toprule
Supporting development & Role and location\\
\midrule
Tseitin encodings, Davis--Putnam elimination, resolution, Horn SAT, and 2-SAT
& Explicit clause certificates, witness reconstruction, least models, and
finite graph algorithms in Section~\ref{ca-solvers}. The original
Davis--Putnam and Tseitin papers provide attribution
\cite{DP1960,Tseitin1968}.\\[4pt]
Forgetting, uniform and Craig interpolation, and finite Beth definability
& Finite Boolean elimination and explicit definitions in
Section~\ref{ca-interpolation}; the proofs use bounded truth tables.\\[4pt]
Finite Heyting algebras, Harrop assumptions, and negative translation
& Semantic and proof-theoretic consequences in
Sections~\ref{sec:ipc-saturation} and~\ref{sec:negative}. Kripke and
Buss--Pudl\'ak provide primary context \cite{kripke1965,buss-pudlak2001}.\\[4pt]
Arithmetic coding and Lean endpoints
& The numerical base and object-derivation plan in
Sections~\ref{sec:base} and~\ref{sec:lean}, connected to the inspected
ProveIt interfaces \cite{proveit}.\\[4pt]
Limits of infinite compactness and intended interpretations
& Section~\ref{sec:limits} separates finite certificates from infinite
paths and from stronger intended-interpretation claims
\cite{burgess1981}.\\
\bottomrule
\end{tabularx}
\end{center}

Cook's paper is used for attribution of the computation-to-propositional
encoding underlying the complexity chapter \cite{cook1971}. The report gives
its own fixed-verifier construction and its own arithmetic audit. Similarly,
no polynomial disjunction-extraction bound from Buss--Pudl\'ak is silently
imported into the report: the bound proved here is primitive recursive.

\subsection{Reading and reproducing the archive}

Read the scope statement and arithmetic base first, then the relevant
mathematical chapter. Section~\ref{sec:ipc-saturation} is self-contained
apart from the common syntax and finite-data conventions. For formalization,
read Section~\ref{sec:lean} together with the supplied Lean plan before
treating any ambient Lean statement as an object-level iPRA theorem.

\begin{center}\small
\begin{tabularx}{\textwidth}{@{}>{\raggedright\arraybackslash}p{.33\textwidth}>{\raggedright\arraybackslash}X@{}}
\toprule
Archive component & Intended use\\
\midrule
\texttt{propositional\_logic\_ipra.pdf} & The complete human-readable argument, bounds,
scope qualifications, source map, and theorem inventory.\\[4pt]
\texttt{report.tex}; \texttt{sections/} & Editable LaTeX source. The
section files are inputs to the main document, not separate standalone
documents. The bibliography is included directly and needs no external
BibTeX database.\\[4pt]
\texttt{code/README.md}; \texttt{code/proplogic.py}
& Executable finite constructions and the data conventions for formulas,
clauses, proof DAGs, and Kripke models.\\[4pt]
\texttt{code/validate.py};
\texttt{code/validation\_results.json}
& The validation driver and its recorded small-instance results.
Run the driver from the code directory with Python 3.10 or later.\\[4pt]
\texttt{code/demo.py};
\texttt{code/demo\_certificates.json}
& Worked certificate outputs for inspection and replay.\\[4pt]
\texttt{code/VERIFICATION.md}
& The meaning and limits of the executable checks. These tests are not
universal mathematical proofs.\\[4pt]
\texttt{lean/Propositional.lean};
\texttt{lean/ObjectExamples.lean}
& Readable Lean specifications and proof candidates, including examples
targeting the inspected object-theory derivation interface.\\[4pt]
\texttt{lean/FormalizationPlan.md}; \texttt{lean/README.md}
& The remaining coding, primitive-recursion, and object-derivation
obligations, together with integration and checking instructions.\\[4pt]
\path{artifacts/repository_audit.json}
& Pinned paths, blob identifiers, and source links for the inspected
upstream integration interfaces. This audit does not constitute a rebuild
of the upstream repository or verification of all its dependencies.\\
\bottomrule
\end{tabularx}
\end{center}

\subsection{Source audit limits}

The repository reference is pinned to the commit recorded in
Section~\ref{sec:lean}; live Wikipedia links are identified by title and
inspection date, not represented as immutable archived revisions. The source
map separates what was inspected from what was proved. Primary papers are
cited for attribution and for distinctions of scope; the report does not
claim those papers already prove the precise iPRA formulations presented
here.

The validation files state the actual computational evidence. In particular,
the absence of a completed Lean build is not concealed by the presence of
Lean syntax, and the checked Python examples are not advertised as proofs
of the universal checker contracts. The full human arguments and explicit
primitive recursive constructions are the mathematical deliverable; the
formalization files specify a clear route to stronger, kernel-checked
evidence.
